\textbf{Phần 1. Trắc nghiệm lựa chọn}
\setcounterpageref{ex}{0}

\begin{ex} %[5D7N36]
Tỉ số của hai số $3$ và $7$ được viết là
\choice
{$\dfrac{7}{3}$}
{\True $\dfrac{3}{7}$}
{$3 \times 7$}
{$7 - 3$}
\loigiai{
Tỉ số của số $3$ và số $7$ là $3 : 7$ hay $\dfrac{3}{7}$.\\
\textbf{Đáp án của câu hỏi là B}
}
\end{ex}

\begin{ex} %[5D7N36]
Phân số $\dfrac{45}{100}$ được viết dưới dạng tỉ số phần trăm là
\choice
{$4{,}5\%$}
{$450\%$}
{\True $45\%$}
{$0{,}45\%$}
\loigiai{
Ta có: $\dfrac{45}{100} = 45\%$.\\
\textbf{Đáp án của câu hỏi là C}
}
\end{ex}

\begin{ex} %[5D7N36]
Số thập phân $0{,}68$ viết dưới dạng tỉ số phần trăm là
\choice
{$6{,}8\%$}
{$0{,}68\%$}
{$680\%$}
{\True $68\%$}
\loigiai{
Ta có: $0{,}68 = 0{,}68 \times 100\% = 68\%$.\\
\textbf{Đáp án của câu hỏi là D}
}
\end{ex}

\begin{ex} %[5D7H36]
Tỉ số phần trăm $25\%$ được viết dưới dạng phân số tối giản là
\choice
{$\dfrac{25}{100}$}
{$\dfrac{1}{5}$}
{\True $\dfrac{1}{4}$}
{$\dfrac{1}{2}$}
\loigiai{
Ta có: $25\% = \dfrac{25}{100} = \dfrac{1}{4}$.\\
\textbf{Đáp án của câu hỏi là C}
}
\end{ex}

\begin{ex} %[5D7H36]
Trong một hộp có $15$ viên bi xanh và $20$ viên bi đỏ. Tỉ số giữa số bi xanh và số bi đỏ là
\choice
{$\dfrac{4}{3}$}
{\True $\dfrac{3}{4}$}
{$\dfrac{3}{7}$}
{$\dfrac{4}{7}$}
\loigiai{
Tỉ số giữa số bi xanh và số bi đỏ là $\dfrac{15}{20} = \dfrac{3}{4}$.\\
\textbf{Đáp án của câu hỏi là B}
}
\end{ex}

\begin{ex} %[5D7H36]
Lớp 5A có $32$ học sinh, trong đó có $18$ bạn nữ. Tỉ số giữa số học sinh nam và số học sinh cả lớp là
\choice
{\True $\dfrac{7}{16}$}
{$\dfrac{9}{16}$}
{$\dfrac{7}{9}$}
{$\dfrac{9}{7}$}
\loigiai{
Số học sinh nam là: $32 - 18 = 14$ (học sinh).\\
Tỉ số của số học sinh nam và số học sinh cả lớp là $\dfrac{14}{32} = \dfrac{7}{16}$.\\
\textbf{Đáp án của câu hỏi là A}
}
\end{ex}

\begin{ex} %[5D7H36]
Phân số $\dfrac{3}{5}$ viết dưới dạng tỉ số phần trăm là
\choice
{$35\%$}
{$30\%$}
{$53\%$}
{\True $60\%$}
\loigiai{
Ta có: $\dfrac{3}{5} = \dfrac{60}{100} = 60\%$.\\
\textbf{Đáp án của câu hỏi là D}
}
\end{ex}

\begin{ex} %[5D7H36]
Tỉ số của hai đại lượng $15\text{ cm}$ và $2\text{ dm}$ là
\choice
{$\dfrac{15}{2}$}
{\True $\dfrac{3}{4}$}
{$\dfrac{4}{3}$}
{$\dfrac{2}{15}$}
\loigiai{
Đổi $2\text{ dm} = 20\text{ cm}$.\\
Tỉ số là: $\dfrac{15}{20} = \dfrac{3}{4}$.\\
\textbf{Đáp án của câu hỏi là B}
}
\end{ex}

\begin{ex} %[5D7V36]
Đoạn thẳng $AB$ dài $12\text{ cm}$, đoạn thẳng $CD$ dài $18\text{ cm}$. Tỉ số độ dài của đoạn thẳng $CD$ so với tổng độ dài hai đoạn thẳng là
\choice
{\True $\dfrac{3}{5}$}
{$\dfrac{2}{3}$}
{$\dfrac{2}{5}$}
{$\dfrac{3}{2}$}
\loigiai{
Tổng độ dài hai đoạn thẳng là: $12 + 18 = 30\text{ (cm)}$.\\
Tỉ số độ dài của đoạn thẳng $CD$ so với tổng độ dài hai đoạn là: $\dfrac{18}{30} = \dfrac{3}{5}$.\\
\textbf{Đáp án của câu hỏi là A}
}
\end{ex}

\begin{ex} %[5D7V36]
Một khu vườn trồng cam và xoài, biết số cây cam bằng $\dfrac{5}{8}$ số cây xoài. Tỉ số giữa số cây xoài và tổng số cây trong vườn là
\choice
{$\dfrac{5}{13}$}
{$\dfrac{8}{5}$}
{\True $\dfrac{8}{13}$}
{$\dfrac{5}{8}$}
\loigiai{
Coi số cây cam là $5$ phần thì số cây xoài là $8$ phần.\\
Tổng số cây trong vườn là $5 + 8 = 13$ phần.\\
Tỉ số giữa số cây xoài và tổng số cây trong vườn là $\dfrac{8}{13}$.\\
\textbf{Đáp án của câu hỏi là C}
}
\end{ex}

\begin{ex} %[5D7V36]
Viết phân số $\dfrac{7}{8}$ dưới dạng tỉ số phần trăm ta được
\choice
{\True $87{,}5\%$}
{$87\%$}
{$78\%$}
{$8{,}75\%$}
\loigiai{
Ta có: $7 : 8 = 0{,}875 = 87{,}5\%$.\\
\textbf{Đáp án của câu hỏi là A}
}
\end{ex}

\begin{ex} %[5D7C36]
Một đàn gà có số gà mái gấp $4$ lần số gà trống. Sau khi bán đi $6$ con gà mái thì số gà mái gấp $3$ lần số gà trống. Tổng số gà ban đầu của đàn là
\choice
{$24$ con}
{$36$ con}
{$40$ con}
{\True $30$ con}
\loigiai{
Vì số gà trống không đổi, coi số gà trống là $1$ phần.\\
Ban đầu số gà mái là $4$ phần. Sau khi bán $6$ con gà mái, số gà mái còn $3$ phần.\\
Do đó $1$ phần ứng với $6$ con gà trống.\\
Số gà trống ban đầu là $6$ con, số gà mái ban đầu là $6 \times 4 = 24$ con.\\
Tổng số gà ban đầu: $6 + 24 = 30$ con.\\
\textbf{Đáp án của câu hỏi là D}
}
\end{ex}

\textbf{Phần 2. Trắc nghiệm đúng sai}
\setcounterpageref{ex}{0}

\begin{ex} %[5D7H36]
Cho biết tỉ số học sinh giỏi và khá của một lớp là $\dfrac{3}{5}$. Xét tính đúng sai của các khẳng định sau:
\choiceTF
{\True Số học sinh giỏi chiếm $3$ phần, số học sinh khá chiếm $5$ phần bằng nhau.}
{Tỉ số phần trăm của số học sinh giỏi so với số học sinh khá là $35\%$.}
{\True Số học sinh giỏi bằng $60\%$ số học sinh khá.}
{Tỉ số của số học sinh khá và tổng số học sinh giỏi, khá là $\dfrac{5}{3}$.}
\loigiai{
\begin{itemize}
    \item Tỉ số $\dfrac{3}{5}$ nghĩa là số học sinh giỏi là $3$ phần, học sinh khá là $5$ phần.
    \item Tỉ số phần trăm: $\dfrac{3}{5} = \dfrac{60}{100} = 60\% \neq 35\%$.
    \item Khẳng định số học sinh giỏi bằng $60\%$ số học sinh khá là đúng.
    \item Tổng số phần là $3 + 5 = 8$. Tỉ số số học sinh khá so với tổng số học sinh giỏi và khá là $\dfrac{5}{8}$ chứ không phải $\dfrac{5}{3}$.
\end{itemize}
\begin{itemize}
    \item a) \textbf{Đúng}.
    \item b) \textbf{Sai}.
    \item c) \textbf{Đúng}.
    \item d) \textbf{Sai}.
\end{itemize}
}
\end{ex}

\begin{ex} %[5D7V36]
Một thư viện có $1\,200$ cuốn sách, trong đó có $480$ cuốn sách khoa học, còn lại là truyện thiếu nhi. Xét tính đúng sai của các khẳng định sau:
\choiceTF
{\True Số truyện thiếu nhi trong thư viện là $720$ cuốn.}
{Tỉ số giữa số sách khoa học và tổng số sách của thư viện là $\dfrac{3}{5}$.}
{\True Tỉ số phần trăm giữa số sách khoa học và tổng số sách là $40\%$.}
{\True Tỉ số giữa số sách khoa học và số truyện thiếu nhi là $\dfrac{2}{3}$.}
\loigiai{
\begin{itemize}
    \item Số truyện thiếu nhi: $1\,200 - 480 = 720$ (cuốn).
    \item Tỉ số sách khoa học và tổng số sách: $\dfrac{480}{1\,200} = \dfrac{2}{5} \neq \dfrac{3}{5}$.
    \item Tỉ số phần trăm: $\dfrac{2}{5} = 40\%$.
    \item Tỉ số sách khoa học và truyện thiếu nhi: $\dfrac{480}{720} = \dfrac{2}{3}$.
\end{itemize}
\begin{itemize}
    \item a) \textbf{Đúng}.
    \item b) \textbf{Sai}.
    \item c) \textbf{Đúng}.
    \item d) \textbf{Đúng}.
\end{itemize}
}
\end{ex}

\textbf{Phần 3. Trả lời ngắn}
\setcounterpageref{ex}{0}

\begin{ex} %[5D7H36]
Viết tỉ số phần trăm $72\%$ dưới dạng phân số tối giản $\dfrac{a}{b}$. Tính giá trị của $a + b$.
\shortans{43}
\loigiai{
\textbf{Đáp số:} $43$\\
$72\% = \dfrac{72}{100} = \dfrac{18}{25} \Rightarrow a = 18, b = 25 \Rightarrow a + b = 18 + 25 = 43$.
}
\end{ex}

\begin{ex} %[5D7H36]
Viết phân số $\dfrac{13}{25}$ dưới dạng tỉ số phần trăm (chỉ ghi số, không ghi kí hiệu $\%$).
\shortans{52}
\loigiai{
\textbf{Đáp số:} $52$\\
$\dfrac{13}{25} = \dfrac{52}{100} = 52\%$.
}
\end{ex}

\begin{ex} %[5D7H36]
Một sợi dây dài $150\text{ cm}$, người ta cắt đi một đoạn dài $60\text{ cm}$. Tỉ số phần trăm của đoạn dây bị cắt đi so với sợi dây ban đầu là bao nhiêu phần trăm? (Chỉ ghi số).
\shortans{40}
\loigiai{
\textbf{Đáp số:} $40$\\
Tỉ số phần trăm: $\dfrac{60}{150} = \dfrac{2}{5} = 40\%$.
}
\end{ex}

\begin{ex} %[5D7V36]
Một cửa hàng có $240\text{ kg}$ gạo tẻ và $160\text{ kg}$ gạo nếp. Tỉ số phần trăm giữa khối lượng gạo nếp và tổng khối lượng gạo của cửa hàng là bao nhiêu phần trăm? (Chỉ ghi số).
\shortans{40}
\loigiai{
\textbf{Đáp số:} $40$\\
Tổng khối lượng: $240 + 160 = 400\text{ kg}$. Tỉ số phần trăm gạo nếp: $160 : 400 = 40\%$.
}
\end{ex}

\begin{ex} %[5D7V36]
Tỉ số của hai số $A$ và $B$ là $\dfrac{3}{8}$. Biết số $A$ là $45$. Tìm số $B$.
\shortans{120}
\loigiai{
\textbf{Đáp số:} $120$\\
Số $B = 45 : 3 \times 8 = 120$.
}
\end{ex}

\begin{ex} %[5D7C36]
Một mảnh đất hình chữ nhật có chu vi $120\text{ m}$. Biết chiều rộng bằng $\dfrac{2}{3}$ chiều dài. Người ta dùng $30\%$ diện tích đất để làm nhà. Tính diện tích phần đất làm nhà theo mét vuông.
\shortans{259,2}
\loigiai{
\textbf{Đáp số:} $259{,}2$\\
Nửa chu vi: $120 : 2 = 60\text{ m}$.\\
Chiều dài: $60 : (2 + 3) \times 3 = 36\text{ m}$.\\
Chiều rộng: $60 - 36 = 24\text{ m}$.\\
Diện tích mảnh đất: $36 \times 24 = 864\text{ m}^2$.\\
Diện tích làm nhà: $864 \times 30\% = 259{,}2\text{ m}^2$.
}
\end{ex}
